\section{Distance under arbitrary binary changes}
\label{sec:distance}

We first prove that fewer than $m^2$ binary changes cannot eliminate
all copies of $H$. We choose one path through the tree on each axis,
with the same node indices, so that its root entries and sampled anchor
and dummy incidences are unchanged. All other entries between variable
positions may change freely. If the edited matrix were $H$-free, its
mode restrictions would then force opposite values on the shared leaf
entry.

\begin{proposition}\label{prop:distance}
Let $h\geq 1$, let $m=2^h$, and let $A=A_h$ be the matrix constructed
above, of order $n=d_hm$, where $d_h=386h+2$.  Every $H$-free binary
$n\times n$ matrix differs from $A$ in at least $m^2$ entries.
Consequently $A$ is $d_h^{-2}$-far from being $H$-free.
\end{proposition}

\begin{proof}
Fix an arbitrary binary matrix $B$ differing from $A$ on a set $E$ of
$M<m^2$ cells.  We will find a copy of $H$ in $B$.

\paragraph{Choosing positions while avoiding the relevant changes.}
On each axis, partition the positions into the anchor groups, the dummy
group, the variable classes at levels $0,\ldots,h-1$ for each sign, and
the one shared leaf class.  All these classes have size $m$.  In
particular, the leaf class occurs once in this partition, despite its
two sign names.

Choose $z$ uniformly from $\{0,\ldots,m-1\}$ and put
\[
 a_i=\left\lfloor\frac{z}{2^{h-i}}\right\rfloor
 \qquad(0\leq i\leq h).
\]
Thus $a_i$ indexes the ancestor at level $i$ of leaf $z$, and
$\lfloor a_i/2\rfloor=a_{i-1}$.  Choose positions
\[
 x_i^\sigma\in R_i^\sigma(a_i),\qquad
 y_i^\sigma\in C_i^\sigma(a_i)
 \qquad(\sigma\in\{+,-\}),
\]
uniformly within their blocks, using a single shared choice on each axis
at the leaf:
$x_h^+=x_h^-$ and $y_h^+=y_h^-$.  Also choose one uniform position from
every anchor group and from each dummy group; denote the dummy choices
by $x_*$ and $y_*$.  Conditional on $z$, all choices in distinct groups
are independent, including choices on opposite axes.

Each variable choice is uniform on its size-$m$ class.  Indeed, a fixed
node at depth $i$ is selected with probability $2^{-i}$, and its block
has size $m/2^i$, so each position in that class is selected with
probability $1/m$.  An anchor or dummy choice is independent of every
choice on the opposite axis.  The root choices are also independent
uniform choices in their fixed size-$m$ blocks.  It follows that, for
every cell $(r,c)$ having at least one anchor or dummy position, or
having both positions at the root level,
\[
 \Pr(r\text{ and }c\text{ are selected})=\frac1{m^2}.
\]
There is only one choice in each size-$m$ class, so each cell occurs in
this calculation once.  Summing over the changed cells of these types
gives
\begin{equation}\label{eq:cellwise-bound}
 \Pr\bigl(\text{some selected cell of these types belongs to }E\bigr)
 \leq\frac{|E|}{m^2}<1.
\end{equation}
Fix a selection avoiding all such changes.  Thus all selected anchor
and dummy incidences, and the selected root entries, agree in $A$ and
$B$.  Changes at every other cell remain unrestricted.

\paragraph{What a mode forbids.}
For each mode $t$, its selected anchor positions form the matrix $S$
in $B$.  They precede all selected non-anchor positions, and their
entries to the selected positions in $\mathsf R_j(t)$ and
$\mathsf C_j(t)$ are the prescribed vectors $e_j$.
Therefore any ordered pair of selected rows in
$\mathsf R_1(t),\mathsf R_2(t)$ and ordered pair of selected columns in
$\mathsf C_1(t),\mathsf C_2(t)$ whose entries in $B$ form
\[
 P=\begin{pmatrix}1&0\\1&1\end{pmatrix}
\]
complete these anchors to a copy of $H$.
Suppose, for a contradiction, that $B$ is $H$-free.  Then this body
matrix cannot occur for any mode.

\paragraph{Propagating the root values.}
The unchanged root entries give
\begin{equation}\label{eq:root-bits}
 B(x_0^+,y_0^+)=1,\qquad B(x_0^-,y_0^-)=0.
\end{equation}
We claim that the same two equalities hold at every level:
\begin{equation}\label{eq:propagation-bits}
 B(x_i^+,y_i^+)=1,\qquad B(x_i^-,y_i^-)=0
 \qquad(0\leq i\leq h).
\end{equation}
Assume they hold at level $i-1$, where $1\leq i\leq h$.
For $V_i^+$ use rows $x_i^+,x_{i-1}^+$ and columns
$y_*,y_{i-1}^+$.  For $V_i^-$ use rows
$x_{i-1}^-,x_i^-$ and columns $y_*,y_{i-1}^-$.  These are in their
designated sets and in the required order because the chosen nodes are
parent and child.  Their body matrices in $B$ are respectively
\[
 \begin{pmatrix}
  1&B(x_i^+,y_{i-1}^+)\\1&1
 \end{pmatrix},\qquad
 \begin{pmatrix}
  1&0\\1&B(x_i^-,y_{i-1}^-)
 \end{pmatrix}.
\]
The dummy incidences supply the left column of ones.  Avoiding $P$
forces
\[
 B(x_i^+,y_{i-1}^+)=1,\qquad
 B(x_i^-,y_{i-1}^-)=0.
\]
To carry these values to level-$i$ columns, put $d=a_i\bmod 2$.
For $W_{i,d}^+$ use rows $x_i^+,x_*$ and columns
$y_{i-1}^+,y_i^+$; for $W_{i,d}^-$ use rows $x_i^-,x_*$ and columns
$y_i^-,y_{i-1}^-$.  Again these positions have the required order
and membership.  The resulting matrices are
\[
 \begin{pmatrix}
  1&B(x_i^+,y_i^+)\\1&1
 \end{pmatrix},\qquad
 \begin{pmatrix}
  B(x_i^-,y_i^-)&0\\1&1
 \end{pmatrix}.
\]
Here the bottom rows are unchanged dummy incidences.  Avoiding $P$
forces exactly \eqref{eq:propagation-bits} at level $i$, completing
the induction.

At level $h$, the two asserted values in
\eqref{eq:propagation-bits} concern the same cell, since the leaf
positions were shared.  This contradiction proves that $B$ contains
$H$.  Since every matrix with fewer than $m^2$ changes has this
property and $n=d_hm$, the normalized distance is at least
$m^2/n^2=d_h^{-2}$.
\end{proof}
